\section{On Animation}\label{appdx: on animation}
We recall the construction of animation following \cite[\S 5.5.8]{HTT} and as discussed in \Cref{sec: lecture 5} and the proof of \Cref{thm: Wagner no-go}.



The construction of animation, or non-Abelian derived functors, dates back to the construction of the cotangent complex, following Quillen and Illusie. We first recall the Dold-Kan correspondence. 
\begin{theorem}[Dold-Kan Correspondence; {\cite[\href{https://stacks.math.columbia.edu/tag/019G}{Tag 019G}]{stacks-project}}]\label{thm: Dold-Kan}
    Let $\Asf$ be an Abelian (1-)category. There is an equivalence of categories between simplicial objects of $\Asf$ and the subcategory of $\Asf$-chain complexes concentrated in non-negative degree. 
\end{theorem}
We now define the cotangent complex (cf. \cite[\href{https://stacks.math.columbia.edu/tag/08PL}{Tag 08PL}]{stacks-project}). 
\begin{definition}[Cotangent Complex]\label{def: cotangent complex}
    The cotangent complex is the chain complex associated to the simplicial $R$-module 
    $$% https://q.uiver.app/#q=WzAsNCxbMywwLCJcXFpaW1JdIl0sWzIsMCwiXFxaWlxcbGVmdFtcXFpaW1JdXFxyaWdodF0iXSxbMSwwLCJcXFpaXFxsZWZ0W1xcWlpcXGxlZnRbXFxaWltSXVxccmlnaHRdXFxyaWdodF0iXSxbMCwwLCJcXGRvdHMiXSxbMywyLCIiLDEseyJvZmZzZXQiOjN9XSxbMywyLCIiLDEseyJvZmZzZXQiOi0zfV0sWzMsMiwiIiwxLHsib2Zmc2V0IjoxfV0sWzMsMiwiIiwxLHsib2Zmc2V0IjotMX1dLFsyLDFdLFsyLDEsIiIsMSx7Im9mZnNldCI6LTN9XSxbMiwxLCIiLDEseyJvZmZzZXQiOjN9XSxbMSwwLCIiLDEseyJvZmZzZXQiOi0xfV0sWzEsMCwiIiwxLHsib2Zmc2V0IjoxfV1d
    \begin{tikzcd}
        \dots & {\ZZ\left[\ZZ\left[\ZZ[R]\right]\right]} & {\ZZ\left[\ZZ[R]\right]} & {\ZZ[R]}
        \arrow[shift right=3, from=1-1, to=1-2]
        \arrow[shift left=3, from=1-1, to=1-2]
        \arrow[shift right, from=1-1, to=1-2]
        \arrow[shift left, from=1-1, to=1-2]
        \arrow[from=1-2, to=1-3]
        \arrow[shift left=3, from=1-2, to=1-3]
        \arrow[shift right=3, from=1-2, to=1-3]
        \arrow[shift left, from=1-3, to=1-4]
        \arrow[shift right, from=1-3, to=1-4]
    \end{tikzcd}$$
    with degeneracy maps augumentation maps under the Dold-Kan correspondence. 
\end{definition}
This construction is motivated by the fact that on smooth (in particular polynomial) $\ZZ$-algebras $R$, the cotangent complex is given by $\LL_{R/\ZZ}\cong\Omega^{1}_{R/\ZZ}[0]$ which has an explicit description in terms of generators and relations. 

More generally, the analogy between derived functors and their non-Abelian counterpart arises from the fact that in both situations one passes to ``free resolutions,'' applies the functor there, and shows the resulting construction is independent of the choices made. 

Before defining animation, we make the following recollections. 
\begin{definition}[Compact Projective Objects]\label{def: compact projective objects}
    Let $\Csf$ be a category with 1-sifted colimits. An object $X\in\Csf$ is compact projective if $\Hom_{\Csf}(X,-)$ commutes with 1-sifted colimits, or equivalently, with filtered colimits and reflexive coequalizers. 
\end{definition}
\begin{example}\label{ex: compact projectives}
    We consider some examples of categories and their compact projective objects. 
    {\begin{center}
        \begin{tabular}{| c | c |}
        \hline 
        $\Csf$ & $\Csf^{\mathsf{cp}}$ \\\hline
        $\Sets$ & finite sets \\\hline
        $\mathsf{CRing}$ & finitely generated polynomial algebras over $\ZZ$ and retracts thereof \\\hline
        $\AbGrp$ & free finitely generated Abelian groups \\\hline
        $\Grp$ & free groups on finitely many generators\\\hline
    \end{tabular}
    \end{center}}
\end{example}
Denoting the full subcategory of $\Csf$ spanned by the compact projective objects $\Csf^{\mathsf{cp}}$, we have by the Yoneda embedding fully faithful embeddings 
$$\Csf^{\mathsf{cp}}\hookrightarrow\mathsf{sInd}^{1}(\Csf^{\mathsf{cp}})\hookrightarrow\Csf$$
where $\mathsf{sInd}^{1}(-)$ is the closure under 1-sifted colimits. 
\begin{definition}[Category Generated by Compact Projectives]\label{def: CPG category}
    Let $\Csf$ be a category with all colimits. $\Csf$ is generated by compact projective objects if there is an equivalence of categories $\mathsf{sInd}^{1}(\Csf^{\mathsf{cp}})\xrightarrow{\sim}\Csf$. 
\end{definition}
We are now ready to define animation of a category in earnest. 
\begin{definition}[Animation of a Category]\label{def: animation of a category}
    Let $\Csf$ be a category admitting all colimits. The animation $\Ani(\Csf)$ of $\Csf$ is the $\infty$-category $\mathsf{sInd}(\Csf^{\mathsf{cp}})$, the closure of (the nerve of) $\Csf^{\mathsf{cp}}$ under $\infty$-categorical sifted colimits. 
\end{definition}
In other words, $\Ani(\Csf)$ is a cocomplete $\infty$-category generated under filtered colimits and geometric realizations of simplicial objects by $\Csf^{\mathsf{cp}}$, or as the localization of the category of simplicial objects of $\Csf$ at the weak equivalences. The animation $\Ani(\Csf)$ of $\Csf$ possesses the universal property that it is initial object in the (non-full) subcategory of cocomplete $\infty$-categories under (the nerve of) $\Csf^{\mathsf{cp}}$ with sifted-colimit preserving functors \cite[\S 5.1.4]{PurityFlatCohomology}. 
\begin{example}\label{ex: animations}
    We consider some examples of categories and their animations.\marginpar{The instructor states that the plural of anima is anima, not animae, despite this being \href{https://logeion.uchicago.edu/morpho/animae}{gramatically incorrect}. The author, having previously studied a classical language, remains insistent on adhering to the correct grammatical conventions.} 
        \begin{center}
        \begin{tabular}{| c | c |}
        \hline 
        $\Csf$ & $\Ani(\Csf)$ \\\hline
        $\Sets$ & $\Ani$, the category of animae or $\infty$-groupoids \\\hline
        $\mathsf{CRing}$ & $\Ani(\Ring)$ the category of animated (commutative) rings \\\hline
        $\AbGrp$ & complexes of Abelian groups conc. in deg. $\geq0$, ie. $\Dscr_{\geq0}(\ZZ)$ \\\hline
        $\Asf$ any Abelian cat. & $\Dscr_{\geq0}(\Asf)$, once again by Dold-Kan \\\hline
    \end{tabular}
    \end{center}
\end{example}
The construction of animation is furthermore functorial, and we can form the animation of functors. For $F:\Csf\to\Dsf$ a 1-sifted colimit preserving functor between cocomplete categories generated by compact projectives, restriction yields a functor $F|_{\Csf^{\mathsf{cp}}}:\Csf^{\mathsf{cp}}\to\Dsf$ and thus a functor to $\Ani(\Dsf)$ by composition. This yields a solid diagram 
$$% https://q.uiver.app/#q=WzAsNCxbMCwwLCJcXENzZl57XFxtYXRoc2Z7Y3B9fSJdLFswLDEsIlxcRHNmIl0sWzIsMSwiXFxBbmkoXFxEc2YpIl0sWzIsMCwiXFxBbmkoXFxDc2YpIl0sWzAsMSwiRnxfe1xcQ3NmXntcXG1hdGhzZntjcH19fSIsMl0sWzEsMiwiXFxBbmkoLSkiLDJdLFswLDMsIlxcQW5pKC0pIiwwLHsic3R5bGUiOnsiYm9keSI6eyJuYW1lIjoiZGFzaGVkIn19fV0sWzMsMiwiXFxBbmkoRikiLDAseyJzdHlsZSI6eyJib2R5Ijp7Im5hbWUiOiJkYXNoZWQifX19XSxbMCwyXV0=
\begin{tikzcd}
	{\Csf^{\mathsf{cp}}} && {\Ani(\Csf)} \\
	\Dsf && {\Ani(\Dsf).}
	\arrow["{\Ani(-)}", dashed, from=1-1, to=1-3]
	\arrow["{F|_{\Csf^{\mathsf{cp}}}}"', from=1-1, to=2-1]
	\arrow[from=1-1, to=2-3]
	\arrow["{\Ani(F)}", dashed, from=1-3, to=2-3]
	\arrow["{\Ani(-)}"', from=2-1, to=2-3]
\end{tikzcd}$$
Since the diagonal composition map $\Csf^{\mathsf{cp}}\to\Ani(\Dsf)$ preserves (nerves of) 1-sifted colimits -- the preservation of $\infty$-categorical sifted colimits is vaccuous as the source is (the nerve of) a 1-category -- the dotted factorization exists by the universal property of animation. 
\begin{definition}[Animation of a Functor]\label{def: animation of a functor}
    Let $F:\Csf\to\Dsf$ be a 1-sifted colimit preserving functor between cocomplete categories generated by compact projectives. The animation $\Ani(F)$ of $F$ is the unique sifted colimit preserving functor $\Ani(\Csf)\to\Ani(\Dsf)$ factoring $\Csf^{\mathsf{cp}}\to\Ani(\Dsf)$. 
\end{definition}
\begin{remark}
    Just as in the case of ordinary categories, animated functors do not necessarily compose well. See \cite[Prop. 5.15]{PurityFlatCohomology} for the requisite hypotheses. 
\end{remark}
Let us return to the cotangent complex. 
\begin{example}
    Let $R$ be a ring and $\Omega^{1}_{(-)/\ZZ}:\mathsf{CRing}\to\AbGrp$ the functor taking each commutative ring to its module of K\"{a}hler differentials. The cotagnent complex $\LL_{-/\ZZ}:\Ani(\Ring)\to\Dscr_{\geq0}(\ZZ)$ is the animation of the functor $\Omega^{1}_{(-)/\ZZ}$. 
\end{example}
We conclude with the following result about the animation of rings and polynomial algebras. 
\begin{proposition}\label{prop: animation of rings and poly algebras agree}
    The animation of the inclusion functor of polynomial $\ZZ$-algebras in finitely many variables into the compact projective objects of commutative rings $\iota:\mathsf{Poly}_{\ZZ}\hookrightarrow\mathsf{CRing}^{\mathsf{cp}}$ is an equivalence $\Ani(\iota):\Ani(\mathsf{Poly}_{\ZZ})\to \Ani(\Ring)$. 
\end{proposition}
\begin{proof}
    Every retract of polynomial algebras over $\ZZ$ in finitely many variables is a quotient, and hence a 1-sifted colimit. In particular, the closure of (the nerve of) $\mathsf{Poly}_{\ZZ}$ under sifted $\infty$-categorical colimits coincides with that of the analogous construction for $\mathsf{CRing}^{\mathsf{cp}}$. 
\end{proof}
